\documentclass[preprint,12pt]{elsarticle}
\usepackage[utf8]{inputenc}
\usepackage{grffile}
\usepackage{longtable}
\usepackage{caption}
\usepackage{setspace}

\usepackage{graphicx}
\usepackage{booktabs}
\usepackage{array}
\usepackage{textgreek}
\usepackage{longtable}

\usepackage{tcolorbox}


\usepackage{placeins}
\usepackage{amsmath,amssymb}

\usepackage{threeparttable}
\usepackage{enumitem}
% \usepackage{multirow}
\usepackage{float}


% \usepackage[authoryear]{natbib}

% Compatibility alias (biblatex -> natbib)
\providecommand{\parencite}{\citep}

\usepackage{hyperref}
\hypersetup{
  colorlinks=true,
  linkcolor=blue,
  citecolor=blue,
  urlcolor=blue
}




\begin{document}


% ========= 1) Title page & front matter (BLINDED for Peer Review) =========
\begin{frontmatter}

\journal{Applied Health Economics and Health Policy}

\title{How Humans Value Delayed Protection: Biological Vulnerability and Intertemporal Health Decisions\tnoteref{t1}}
\tnotetext[t1]{Job Market Paper with Professor Dohyeong Kim}

\author{Yichao Jin\fnref{fn1}}
\ead{yichao.jin@utdallas.edu}
\author{Dohyeong Kim\fnref{fn1}}

\address[fn1]{School of Economic, Political and Policy Sciences, The University of Texas at Dallas}

\begin{abstract}
Vaccination timing critically determines the duration of biological exposure to infectious 
risk, yet many individuals delay immunisation even when vaccines are free and accessible. 
While traditional models focus on binary uptake, this study quantifies the behavioral cost 
of waiting as an intertemporal health investment. Using a large-scale discrete choice 
experiment (N=1,027) in Wuhan, China, we recover structural time-preference parameters 
characterizing sensitivity to delayed biological protection.

Results reveal pronounced present-oriented preferences: respondents strongly discount 
delayed vaccination, with hyperbolic parameter $\kappa = 0.225$ and exponential 
$\delta = 0.160$. These parameters translate into an average marginal willingness-to-accept (MWTA) 
of approximately 47 RMB per month of waiting. Delay aversion is systematically patterned by 
biological and social vulnerability—specifically, individuals with lower 
institutional trust exhibit substantially steeper discounting and higher MWTA (62.3 RMB) 
compared to high-trust respondents (28.7 RMB). 

To ensure the robustness of our welfare estimates, we utilize Mixed Logit models to identify individual-level taste heterogeneity and Conditional Logit models for stable population-level MWTA derivation. By linking behavioral parameters to biological risk and uncertainty, we show that 
vaccination delays impose regressive behavioral burdens. The findings provide actionable inputs for designing time-consistent public health 
strategies that minimize the behavioral cost of waiting and improve the efficiency and 
equity of preventive-health delivery.
\end{abstract}

\begin{keyword}
Vaccination timing; Intertemporal preferences; Health capital; Biological vulnerability; Institutional trust; Delay aversion; Discrete choice experiment; SEIR dynamics; Economic policy
\end{keyword}

\end{frontmatter}

\clearpage

\section{Introduction}
\doublespacing
Timely access to biological protection is a central determinant of human exposure to infectious disease and the preservation of health capital \cite{Castioni2024, Shankar2024}. While vaccination coverage is often treated as a binary outcome—vaccinated or not—the timing of vaccination critically shapes the duration for which individuals remain unprotected and biologically vulnerable. Even modest delays can extend periods of exposure and shift outbreak dynamics, making the intertemporal timing of uptake as critical as aggregate population coverage \cite{Rasmussen2021}. Delays in vaccination therefore prolong exposure to infection risk, during which individuals face potential health deterioration and uncertainty about future protection. Understanding how humans evaluate such delays is essential for explaining observed patterns of preventive health behaviour under infectious risk.

From an intertemporal perspective, vaccination timing represents a health investment decision in which immediate costs—such as effort, inconvenience, or perceived side effects—are traded off against delayed and uncertain biological benefits \cite{Frederick2002}. When protection is postponed, individuals must decide whether the future reduction in infection risk sufficiently compensates for remaining exposed in the present. These trade-offs are not uniform across populations. Institutional trust \cite{Krastev2023}, perceived biological vulnerability \cite{Bok2024}, and stable preventive-health routines \cite{Kumar2025} shape how future protection is perceived and discounted, giving rise to systematic heterogeneity in tolerance for delayed protection. These determinants appear across diverse institutional settings, reflecting general features of intertemporal health decision-making.

A growing body of research documents that vaccination timing is influenced by behavioural factors such as impatience, risk perception, and institutional trust, alongside logistical access constraints \cite{Anderson2023, Robertson2024, Etowa2024}. However, much of this literature continues to treat waiting time as an administrative friction rather than as a biologically meaningful interval of continued exposure \cite{Lievre2024, Tran2025, Whitaker2026}. From a physiological perspective, delay is not neutral: each additional period without protection entails ongoing exposure during which infection risk, potential morbidity, and uncertainty accumulate. Recent large-scale evidence highlights that addressing these behavioral drivers is key to improving uptake \cite{Owens2026, Siram2024}. Despite this, few discrete choice experiments have directly quantified how individuals value delayed protection as an intertemporal health trade-off \cite{Yue2021, Gong2020}, and existing studies of health-related discounting have not linked sensitivity to delay to structural determinants such as institutional trust or chronic medical vulnerability \cite{vanderPol2001}. 

To address these gaps, this study examines how individuals value the \emph{timing} of vaccination as an intertemporal health investment. We use a discrete choice experiment (DCE) administered to 1{,}027 adults in Wuhan, China, in which waiting time is embedded directly into vaccine attributes. This design allows us to recover formal time-preference parameters that characterize sensitivity to delayed biological protection. Rather than treating waiting as a purely logistical friction, we conceptualize delay as a biologically meaningful interval during which individuals remain exposed to infection risk and potential health damage. From this perspective, vaccination timing reflects how individuals trade off immediate burdens against the expected preservation of future health under uncertainty.

Wuhan provides a particularly informative empirical setting for studying these processes. As a city with high population density, early pandemic exposure, and strong public-health visibility, Wuhan constitutes a high-salience environment in which the biological consequences of delayed protection are readily understood by respondents. This salience sharpens the behavioral signal linking perceived vulnerability, trust in institutions, and tolerance for waiting, facilitating identification of intertemporal preferences that are otherwise difficult to observe. While absolute valuations of delay may vary across contexts, the underlying mechanisms identified here—risk-weighted time preferences, present-oriented valuation of protection, and heterogeneity by biological and social vulnerability—reflect general features of human health behavior rather than idiosyncratic local conditions. Institutional trust, in particular, has emerged as a critical moderator of vaccination efficiency and equity in diverse settings \cite{Sohns2024, MacDonald2025}. 

This study makes three contributions. First, it provides a DCE-based quantification of intertemporal preferences for vaccination, recovering pronounced short-run impatience toward delayed protection, consistent with present-biased discounting in health investment decisions. Second, it demonstrates that sensitivity to delay is systematically patterned by biological and behavioral characteristics, including chronic medical conditions, preventive-health routines, and institutional trust, indicating that time preferences are shaped by perceived vulnerability and uncertainty rather than being purely individual traits. Third, by translating delay sensitivity into marginal willingness-to-accept measures, the analysis reveals that identical waiting periods impose unequal experiential burdens across the population, contributing to inequalities in the timing of health protection. Finally, by integrating these behavioral parameters into a transmission model, we estimate the social shadow price of delay, echoing recent structural analyses of the combined impact of vaccines and behavior on population health outcomes \cite{Atkeson2024}. 

Together, these findings position vaccination delay as a biologically consequential dimension of preventive behavior and underscore the importance of accounting for intertemporal preferences when studying health-related decision-making. 


% ---------------------------------------------
\section{Theoretical Framework}

Vaccination timing can be conceptualized as an investment in health capital,
in which individuals weigh the immediate utility costs of uptake—such as
effort, inconvenience, or anxiety—against the discounted preservation of
future health stocks \cite{Frederick2002}. From this perspective, vaccination
decisions are inherently intertemporal: postponing protection extends the
period during which individuals remain biologically exposed to infectious
risk. The discount parameters $\kappa$ and $\delta$ therefore reflect not
merely abstract time preferences, but valuations shaped by individuals’
biological vulnerability and institutional environment. When future health
returns are heavily discounted, even short delays impose substantial
behavioural disutility, leading to underinvestment in timely preventive
health.

Discrete choice experiments (DCEs) provide a structured approach to
quantifying these intertemporal valuations by embedding vaccination delay
directly into the latent utility associated with each alternative. Under an
exponential specification, the utility of delayed protection declines at a
constant rate $\delta$, consistent with time-consistent valuation. In
contrast, a quasi-hyperbolic specification allows discounting to be steeper
over short horizons, capturing present-oriented preferences through the
parameter $\kappa$. Rather than serving as purely technical objects, these
parameters summarize how individuals trade off immediate burdens against
delayed biological protection when exposure risk persists during waiting.

To capture how individuals \emph{experience} waiting time rather than simply
its objective duration, we derive a measure of subjective delay that links
formal discounting to perceived temporal burden. This measure reflects the
extent to which waiting for vaccination is psychologically inflated relative
to calendar time, representing a behavioural “time tax” associated with
remaining unprotected. Within this framework, the delay coefficient recovered
from the DCE captures structural sensitivity to ongoing exposure rather than
a purely logistical inconvenience. Variation in this sensitivity is expected
to arise systematically from biological and health-related characteristics—
such as age and chronic medical conditions—as well as from institutional
trust, which stabilizes expectations about future protection and reduces
uncertainty surrounding delayed health returns. For clarity, we refer to the
ratio between perceived and objective waiting time as the \emph{Behavioral
Delay Multiplier (BDM)}. While BDM reflects aggregate behavioral sensitivity, we conceptualize it as a compound measure incorporating both subjective time inflation and perceived institutional reliability.



\section{Methods}

Our methodological approach integrates discrete choice modelling with a behavioral-economic
framework to quantify how individuals value timely vaccination. Rather than treating waiting
time as a logistical inconvenience, we embed delay directly into vaccine attributes and recover
time-discount parameters from respondents’ revealed trade-offs. This approach yields two 
economically interpretable quantities: the exponential discount factor $\delta$, capturing
how the value of protection declines with each month of delay, and the quasi-hyperbolic 
parameter $\kappa$, capturing present bias and the disproportionate weight placed on 
immediacy.

These parameters provide a welfare-relevant measure of the behavioral cost of delay. By 
mapping the delay coefficient onto marginal willingness-to-accept (MWTA), we estimate the 
implicit compensation individuals require to tolerate additional waiting time. Subgroup-specific
estimates allow us to examine how delay aversion varies across demographic, behavioral, and 
institutional characteristics, providing a basis for identifying which groups may benefit most 
from reduced waiting times or early-access incentives.

The main text focuses on the conceptual and policy relevance of these parameters. All formal 
derivations—including the mapping from choice-model coefficients to $\delta$, 
$\kappa$, and MWTA—are provided in Appendix~C. This separation ensures transparency 
while keeping the primary narrative centred on the economic interpretation of timeliness in 
vaccination.

\subsection{Study Design}

We conducted a discrete choice experiment (DCE) to examine behavioral preferences for timely 
vaccination against COVID-like diseases (CLD).\footnote{Respondents were asked to consider a 
future “COVID-like disease” (CLD) scenario, defined as a hypothetical respiratory infection with 
severity and transmissibility comparable to COVID-19. This neutral framing enabled preference 
elicitation without anchoring responses to specific waves, variants, or personal experiences.} 


The CLD framing serves an important methodological function. By presenting a hypothetical 
respiratory pathogen with COVID-like transmissibility and severity, the design avoids anchoring 
responses to particular epidemic waves, personal experiences, or time-specific emotions 
associated with the original pandemic. This reduces recall bias, emotional spillover, and 
context-specific heuristics that may distort valuation of delay during or immediately after an 
actual outbreak. The CLD scenario therefore enables recovery of structural time-preference 
parameters—such as $\delta$, $\kappa$, and MWTA—that generalise beyond COVID-19 and 
apply to a broader class of high-transmission respiratory diseases. This framing strengthens 
external validity while maintaining sufficient realism for respondents to make meaningful 
intertemporal trade-offs.


DCEs enable respondents to evaluate structured trade-offs between vaccine attributes and are widely applied in health-preference research \cite{Louviere2000,Hensher2015}. Although stated preferences may not perfectly mirror revealed behaviour, DCEs are extensively validated in preventive-health contexts and reliably recover marginal valuations of delay, convenience, and risk \cite{Louviere2000,Hensher2015,Johnson2013}.
 The survey was 
administered online in early 2025, when vaccination services were broadly available and respondents 
could reflect on vaccination timing outside acute epidemic pressures.

\subsection{Sampling and Participants}

The survey targeted adults living in Wuhan, China, using sampling strata for age, gender, and 
district to approximate the demographic composition of the city’s adult population. Wuhan 
provides a theoretically meaningful setting for studying socially patterned vaccination behavior: 
its early and intense pandemic experience, strong public-health infrastructure, and high 
institutional visibility create a context in which differences in trust, information, and preventive 
routines are particularly salient. Eligibility criteria required respondents to be aged 18 or older, 
reside in Wuhan, and complete all survey modules. Standard attention and data-quality checks were 
applied. The final analytic sample consisted of 1,027 respondents. A comparison with the 2020 
Wuhan Census benchmarks is provided in Appendix~B.

\subsection{DCE Attributes, Choice Tasks, and Experimental Design}

Vaccine profiles were defined by five attributes: vaccination delay (0--6 months), vaccine efficacy, expected side effects, vaccine origin, and cash incentives. Each respondent completed six choice tasks, with each task presenting two vaccine alternatives and an opt-out option. An efficient fractional-factorial design was employed to ensure attribute balance and identification of main effects.

Prior to the main survey, the instrument was validated through a pilot study ($n=60$) to assess attribute clarity and cognitive burden. Pilot results informed the selection of priors for the efficient design algorithm and indicated a mean completion time of 6.1 minutes, suggesting that task complexity was appropriate for the target population. The full choice card, including attribute definitions and pilot-based validity diagnostics, is provided in Appendix~A.

\subsection{Statistical Analysis}

We estimated mixed logit models to recover preference distributions and account for unobserved heterogeneity in vaccination choices. Attribute coefficients were then used to derive implied exponential ($\delta$) and hyperbolic ($\kappa$) discount parameters, following established behavioral-economic formulations \cite{vanderPol2001}. These parameters quantify the degree to which respondents devalue delayed protection. Subgroup analyses examined variation by demographics, health behaviors, chronic medical conditions, and institutional trust, illuminating the social patterning of delay aversion.

We also computed marginal willingness-to-accept (MWTA) for vaccination delay by taking the ratio of the delay coefficient to the incentive coefficient. MWTA provides an economically interpretable measure of the implicit compensation individuals require to tolerate additional waiting time. All analyses were conducted in R using the \texttt{mlogit} and \texttt{gmnl} packages. Full model specifications, estimation details, and robustness checks—including alternative model structures, centering approaches, and repeated-choice diagnostics—are presented in Appendix~C. Specifically, to address potential collinearity between present bias ($\kappa$) and long-term discounting ($\delta$), we utilize the individual-level heterogeneity recovered via Mixed Logit to ensure that structural parameters are identified from variation in choices across multiple delay horizons (0, 1, 3, and 6 months) rather than a single linear slope.

\clearpage

\section{Results}\label{sec:results}

\subsection{Main Model Estimates}

Table~\ref{tab:models} summarizes the estimates from the conditional, multinomial, and mixed logit models based on the discrete choice experiment. To ensure comparability, coefficients were standardized to represent the change in utility associated with a one–standard-deviation change in each attribute.

In all models, \textit{vaccine delay} was negatively associated with the probability of vaccine uptake, confirming that a longer delay substantially reduces willingness to vaccinate. Respondents showed strong positive preferences for vaccine efficacy and cash incentives, whereas side effects had weaker and less consistent associations with uptake.


Comparing magnitudes illustrates behavioral trade-offs: a one–standard-deviation increase in cash incentive offset nearly the entire disutility of a one–standard-deviation increase in vaccine delay, indicating that monetary incentives can effectively compensate for delay costs. Improvements in vaccine efficacy yielded similar gains, suggesting that respondents balance time, money, and protection when deciding whether to vaccinate.

The mixed logit model provided the best overall fit based on the log-likelihood and Akaike Information Criterion (AIC). Its random-parameter specification captures intra-individual preference correlation, providing a more flexible representation of heterogeneity. The significant standard deviation of the \textit{multiple-delays} parameter in the mixed logit model indicates substantial variation in time sensitivity across respondents. Accordingly, the standardized coefficient for vaccine delays ($-0.353$) from the mixed logit model was used to estimate the hyperbolic and exponential discount parameters reported in the following tables.

Table~\ref{tab:models} shows the comparison of logit models.


\begin{table}[!t]
\centering
\caption{Comparison of logit models (dependent variable: choice)}
\label{tab:models}
\begin{threeparttable}
\begin{tabular}{lccc}
\toprule
 & Conditional logit & MNL (ASCs) & Mixed logit \\
\midrule
Vaccine delays (std) & $-0.131^{***}$ (0.019) & $-0.125^{***}$ (0.020) & $-0.353^{***}$ (0.027) \\
Vaccine efficacy (std) & $0.153^{***}$ (0.036) & $0.163^{***}$ (0.037) & $0.218^{***}$ (0.041) \\
Side effects (std) & $-0.042$ (0.027) & $-0.082^{***}$ (0.030) & $0.008$ (0.031) \\
Cash incentives (std) & $0.327^{***}$ (0.018) & $0.313^{***}$ (0.019) & $0.431^{***}$ (0.025) \\
Vaccine origin (imported) & $0.244^{***}$ (0.038) & $0.239^{***}$ (0.038) & $0.093^{**}$ (0.047) \\
Opt-out ASC &  & $-0.214^{**}$ (0.093) & $-0.287^{**}$ (0.118) \\
ASC: B &  & $-0.108^{***}$ (0.033) &  \\
ASC: C &  & $-0.326^{***}$ (0.100) &  \\
SD vaccine delays (std) &  &  & $1.303^{***}$ (0.073) \\
Num. obs. & 6162 & 6162 & 6162 \\
Log likelihood & -6101.771 & -6096.612 & -5615.522 \\
AIC & 12215.543 & 12207.225 & 11245.044 \\
\bottomrule
\end{tabular}
\begin{tablenotes}[flushleft]
\footnotesize
\item Note: Standard errors in parentheses. $^{*}p<0.1$, $^{**}p<0.05$, $^{***}p<0.01$.
\end{tablenotes}
\end{threeparttable}
\end{table}

\clearpage

\subsection{Discount-Rate Estimation}

From a biological perspective, steeper short-run discounting implies a higher
perceived cost of remaining unprotected during ongoing exposure.
Figure~\ref{fig:figure2} presents the estimated discounting profiles from the exponential and
hyperbolic models, fitted to the observed valuation points using a
sum-of-squared-errors procedure. The \textit{x}-axis shows the delay from 0 to 6 months,
and the \textit{y}-axis shows the corresponding discount factor. As expected, both curves
begin at full value at zero delay and decline monotonically with increasing waiting time.
The exponential model yields an estimated discount factor of $\delta = 0.160$, while the
hyperbolic model yields $\kappa = 0.225$. The hyperbolic curve exhibits a sharp initial
decline, particularly within the first three months, indicating pronounced present-oriented
preferences and a strong valuation of immediacy in health protection. Appendix~D reports
subgroup-specific exponential and hyperbolic fits.

Beyond their behavioural interpretation, these discount parameters have
well-defined economic meaning. In standard intertemporal choice and health-
investment models, discount factors represent the rate at which individuals
trade off present versus future health returns. Accordingly, $\delta$ and $\kappa$ cap-
ture not only impatience but also the degree of intertemporal consistency
and investment orientation in vaccination decisions, linking our estimates to
foundational health-capital models and preventive behaviour theory.
From a human-biology perspective, steeper short-run discounting implies a higher marginal biological cost of remaining unprotected, consistent with elevated perceived susceptibility to infection among medically vulnerable respondents.
To link estimated discounting to experiential welfare, we also derive a Behavioral
Delay Multiplier (BDM) that captures the inflation of subjective waiting
time.

\paragraph{External Benchmarking of Discount Parameters}
To assess external validity, we benchmarked our estimates against international evidence.
Early DCE-based studies in the United Kingdom reported implied discount rates of
0.055--0.091 for one’s own health and 0.078--0.147 for others’ health \cite{vanderPol2001}.
Our estimates ($\delta = 0.160$, $\kappa = 0.225$) indicate steeper short-run
impatience but remain within the broader empirical range of health-related discounting.
Recent work documents similar patterns: Guillon, Béraud, and Vallée \cite{Guillon2024}
identify impatience as a principal driver of COVID-19 vaccination; Attema, Brouwer, and
Claxton \cite{Attema2013} and Halilova et al. \cite{Halilova2025} demonstrate that present bias
and higher discount rates reduce initiation and booster uptake; and Soofi, Rashidi, and
Tavakoli \cite{Soofi2025} estimate comparable $\beta$--$\delta$ parameters for preventive
behaviours in Iran. Taken together, this evidence suggests that the impatience observed in
Wuhan reflects a systematic behavioural regularity rather than a context-specific anomaly.

\paragraph{Contextual Salience of the Wuhan Setting}
The magnitude of short-run impatience may partly reflect the heightened salience of
vaccination during Wuhan’s early-pandemic experience as the initial global epicentre of
COVID-19, characterised by elevated risk awareness and institutional visibility. However,
the underlying mechanisms—present bias, trust-conditioned valuation, and sensitivity to
delay—are consistent with patterns documented across diverse international settings,
suggesting that the behavioural processes identified here are generalisable even if the
Wuhan context accentuates their magnitude.


\begin{figure}[H]
    \centering
    \includegraphics[width=1.0\linewidth]{Overall Discounting Hyoerbolic vs Expondential.png}
    \caption[Observed vs. fitted discounting curves.]{Comparison of observed discounting with hyperbolic and exponential fits.}
    \label{fig:figure2}
    Alt text: A line chart comparing observed discounting behaviour with fitted hyperbolic and exponential curves, showing close alignment in early months and divergence over longer delays.
\end{figure}




\clearpage

\subsection{Subgroup Heterogeneity}

Table~\ref{tab:subgroups_long} summarizes the estimated hyperbolic and exponential discount rates across
demographic, health–behavioral, prior vaccine reaction, and institutional‐trust subgroups.
For each subgroup, $\beta_{\text{delay}}^{(\mathrm{std})}$ represents the standardized delay coefficient, while
$\beta_{\text{delay}}^{(\mathrm{month})}$ denotes the unstandardized marginal utility loss per month of vaccination
delay. The transformation from waiting-time coefficients to the exponential and hyperbolic discount 
parameters ($\delta$, $\kappa$) follows the derivation provided in Appendix~C.6.
Together, these indicators reveal how different population groups value the immediacy
of protection.


\noindent\textbf{Demographic heterogeneity.}
Clear demographic gradients emerged under both discounting models. Women exhibited
slightly higher discount parameters than men, indicating greater impatience toward delayed
vaccination and a stronger preference for prompt protection. Rural respondents,
however, demonstrated lower discount rates relative to urban residents, suggesting stronger
willingness to wait when faced with delays in vaccine availability. Age‐related differences
followed a nonlinear pattern: middle‐aged adults (45–65 years) were the most patient, whereas
younger ($<35$ years) and older ($\geq 65$ years) adults displayed steeper discounting,
consistent with previous evidence that these groups place greater weight on short‐term
benefits in health decisions. Educational differences were similarly nonmonotonic.
Individuals with moderate education levels showed the highest patience, whereas those with
very low or very high educational attainment exhibited stronger present bias. Both hyperbolic
($\kappa$) and exponential ($\delta$) estimates produced low standard errors across subgroups,
supporting the stability and internal consistency of the demographic patterns.

\noindent\textbf{Health--behavioral and institutional heterogeneity.}
Time preferences and MWTA varied meaningfully across behavioral and institutional dimensions.
In our harmonized analysis, respondents are partitioned into subgroups consistent with the unified definitions in Chapter 2.
Institutional trust emerges as a primary driver of delay valuation: respondents with low institutional trust (including neutral and distrust cohorts) require 62.3 RMB per month to tolerate delay, more than double the 28.7 RMB required by high-trust individuals. 

Smoking status and physical activity further capture how routine health behaviors
shape intertemporal vaccine decisions. For smoking status, we classify smokers as individuals who smoke at least sometimes (n=716) and non-smokers as those who have never smoked (n=311).
Contrary to the generalized impulsivity hypothesis, smokers in our sample require an MWTA of approximately 45.5 RMB per month to accept delay, compared to 48.1 RMB among non-smokers.
This suggests that delay aversion in the health domain may be distinct from domain-specific impulsivity or is significantly shaped by the pandemic context. This ``smoking paradox'' --- where smokers exhibit greater patience for vaccination --- may reflect a selection effect among survivors or a context-specific valuation where the relative urgency of vaccination is lower compared to existing physiological stressors.

Methodologically, we utilize Mixed Logit models to identify individual-level taste heterogeneity and Conditional Logit models for stable, population-level welfare estimation. This dual-modeling strategy ensures that while we identify who is sensitive to delay, our monetary valuations remain robust for policy thresholds.
the lowest discount rates within their respective categories, indicating greater willingness to
tolerate delayed vaccination. One interpretation is that individuals with more stable, habitual
health routines may adopt a more flexible temporal orientation toward preventive interventions.
By contrast, respondents with less regular health behaviors displayed steeper short-run
discounting, consistent with behavioral-economic evidence linking irregular routines to
stronger present bias.

Respondents reporting chronic medical conditions exhibited steeper discounting than those
without such conditions. Consistent with a biological-risk interpretation, respondents
reporting chronic medical conditions exhibited steeper discounting, implying that remaining
unprotected carries a higher expected biological cost for individuals with elevated
susceptibility to severe infection. Access to healthcare facilities contributed more modestly to
delay sensitivity: respondents unable to reach a facility within 30 minutes showed slightly
higher discount rates, indicating greater impatience. Prior adverse reactions to vaccines
produced minimal differences, suggesting that past side-effect experiences did not
substantially alter delay valuations in the context of COVID-like disease vaccination.






\clearpage
\begin{longtable}{lcccc}
\caption{Hyperbolic ($\kappa$) vs.\ Exponential ($\delta$) parameters across demographic, behavioral, and institutional‐trust subgroups} \label{tab:subgroups_long} \\
\toprule
\textbf{Group} & \textbf{$n$} & \textbf{$\beta_{\text{delay}}$ (std)} & \textbf{$\beta_{\text{delay}}$ (per mo)} & \textbf{$\kappa$ / $\delta$} \\
\midrule
\endfirsthead

\multicolumn{5}{l}{\textit{Table~\ref{tab:subgroups_long} (continued)}}\\
\toprule
\textbf{Group} & \textbf{$n$} & \textbf{$\beta_{\text{delay}}$ (std)} & \textbf{$\beta_{\text{delay}}$ (per mo)} & \textbf{$\kappa$ / $\delta$} \\
\midrule
\endhead

\bottomrule
\multicolumn{5}{r}{\textit{Continued on next page}}\\
\endfoot

\bottomrule
\multicolumn{5}{p{0.9\textwidth}}{\footnotesize
Notes: Standard errors in parentheses. $n$ indicates the number of respondents per subgroup.
$p$‐values adjusted via Benjamini–Hochberg false‐discovery rate ($q = 0.10$); Holm–Bonferroni
corrections yield similar results. Significance levels: $^{*}p<0.1$; $^{**}p<0.05$; $^{***}p<0.01$.}\\
\endlastfoot

Overall & 1{,}027 & $-0.353^{***}$ (0.027) & $-0.160^{***}$ (0.012) & 0.225 / 0.160 \\
\addlinespace
\multicolumn{5}{l}{\textit{Gender}}\\
Female & 545 & $-0.359^{***}$ (0.034) & $-0.163^{***}$ (0.015) & 0.230 / 0.163 \\
Male & 482 & $-0.346^{***}$ (0.036) & $-0.157^{***}$ (0.016) & 0.219 / 0.157 \\
\addlinespace
\multicolumn{5}{l}{\textit{Residence}}\\
Rural & 555 & $-0.287^{***}$ (0.033) & $-0.130^{***}$ (0.015) & 0.174 / 0.130 \\
Urban & 472 & $-0.433^{***}$ (0.038) & $-0.196^{***}$ (0.017) & 0.290 / 0.196 \\
\addlinespace
\multicolumn{5}{l}{\textit{Age group}}\\
18--34 & 217 & $-0.327^{***}$ (0.054) & $-0.148^{***}$ (0.024) & 0.204 / 0.148 \\
35--44 & 320 & $-0.420^{***}$ (0.042) & $-0.190^{***}$ (0.019) & 0.279 / 0.190 \\
45--54 & 231 & $-0.292^{***}$ (0.051) & $-0.132^{***}$ (0.023) & 0.177 / 0.132 \\
55--65 & 136 & $-0.237^{***}$ (0.059) & $-0.107^{***}$ (0.027) & 0.137 / 0.107 \\
65+ & 123 & $-0.477^{***}$ (0.070) & $-0.216^{***}$ (0.032) & 0.327 / 0.216 \\
\addlinespace
\multicolumn{5}{l}{\textit{Education}}\\
No formal education & 65 & $-0.442^{***}$ (0.093) & $-0.201^{***}$ (0.042) & 0.298 / 0.201 \\
Primary school & 226 & $-0.500^{***}$ (0.053) & $-0.227^{***}$ (0.024) & 0.348 / 0.227 \\
Middle school & 264 & $-0.265^{***}$ (0.049) & $-0.120^{***}$ (0.022) & 0.158 / 0.120 \\
High school & 210 & $-0.277^{***}$ (0.052) & $-0.125^{***}$ (0.024) & 0.166 / 0.125 \\
College/University & 160 & $-0.282^{***}$ (0.058) & $-0.128^{***}$ (0.026) & 0.170 / 0.128 \\
Graduate & 102 & $-0.458^{***}$ (0.063) & $-0.208^{***}$ (0.029) & 0.311 / 0.208 \\
\addlinespace
\multicolumn{5}{l}{\textit{Smoking status}}\\
Never & 311 & $-0.367^{***}$ (0.043) & $-0.167^{***}$ (0.020) & 0.236 / 0.167 \\
Sometimes & 402 & $-0.392^{***}$ (0.039) & $-0.178^{***}$ (0.018) & 0.256 / 0.178 \\
Often & 314 & $-0.290^{***}$ (0.044) & $-0.131^{***}$ (0.020) & 0.176 / 0.131 \\
\addlinespace
\multicolumn{5}{l}{\textit{Physical activity}}\\
Rare & 297 & $-0.455^{***}$ (0.046) & $-0.206^{***}$ (0.021) & 0.308 / 0.206 \\
Sometimes & 439 & $-0.304^{***}$ (0.037) & $-0.138^{***}$ (0.017) & 0.186 / 0.138 \\
Quite often & 291 & $-0.325^{***}$ (0.043) & $-0.147^{***}$ (0.020) & 0.202 / 0.147 \\
\addlinespace
\multicolumn{5}{l}{\textit{Medical condition}}\\
No & 781 & $-0.300^{***}$ (0.029) & $-0.136^{***}$ (0.013) & 0.183 / 0.136 \\
Yes & 119 & $-0.355^{***}$ (0.070) & $-0.161^{**}$ (0.032) & 0.227 / 0.161 \\
\addlinespace
\multicolumn{5}{l}{\textit{Healthcare access}}\\
No access & 460 & $-0.335^{***}$ (0.037) & $-0.152^{***}$ (0.017) & 0.211 / 0.152 \\
Access & 567 & $-0.368^{***}$ (0.034) & $-0.167^{***}$ (0.015) & 0.237 / 0.167 \\
\addlinespace
\multicolumn{5}{l}{\textit{Prior vaccine reaction}}\\
No & 950 & $-0.352^{***}$ (0.027) & $-0.159^{***}$ (0.012) & 0.224 / 0.159 \\
Yes & 77 & $-0.372^{***}$ (0.092) & $-0.168^{**}$ (0.042) & 0.240 / 0.168 \\
\addlinespace
\multicolumn{5}{l}{\textit{Institutional trust}}\\
Strongly distrust & 93 & $-0.327^{***}$ (0.072) & $-0.148^{***}$ (0.033) & 0.204 / 0.148 \\
Somewhat distrust & 133 & $-0.317^{***}$ (0.072) & $-0.144^{***}$ (0.032) & 0.196 / 0.144 \\
Neutral & 261 & $-0.510^{***}$ (0.052) & $-0.231^{***}$ (0.023) & 0.356 / 0.231 \\
Somewhat trust & 267 & $-0.278^{***}$ (0.045) & $-0.126^{***}$ (0.020) & 0.167 / 0.126 \\
Strongly trust & 179 & $-0.268^{***}$ (0.054) & $-0.122^{***}$ (0.025) & 0.160 / 0.122 \\
\end{longtable}

\clearpage



To synthesize the subgroup patterns, Figure~\ref{fig:subgroup_kappa} visualizes the estimated
hyperbolic discount parameters ($\kappa$) across demographic, behavioral, health-status,
and institutional-trust groups. The figure highlights the broad dispersion in delay sensitivity,
with certain groups exhibiting markedly steeper present-oriented preferences than others.
Higher values of $\kappa$ cluster among subgroups characterised by greater uncertainty
or social vulnerability, whereas lower values appear among groups with stronger health
engagement or stability in preventive behaviours. The visualization makes clear that time
preferences for vaccination are heterogeneous rather than uniform, reinforcing the need for
tailored strategies that account for these uneven behavioral costs of waiting.


\begin{figure}[H]
    \centering
    \includegraphics[width=\linewidth]{Subgroup Variation in Hyperbolic Discount Parameters.png}
    \caption{Subgroup variation in hyperbolic discount parameters ($\kappa$). Higher values indicate stronger present-oriented preferences (greater impatience toward vaccination delay). Clear social and behavioral gradients appear across education, age, physical activity, medical vulnerability, and trust in government.}
    \label{fig:subgroup_kappa}
\end{figure}





\clearpage
\subsection{Marginal Willingness-to-Accept (MWTA)}

Table~\ref{tab:mwta_subgroups} reports the implied marginal willingness-to-accept (MWTA) for one 
month of delay. Across the full sample, respondents required approximately 45--56 RMB in 
compensation per month of waiting, reflecting substantial short-run impatience. These MWTA 
estimates align closely with the subgroup differences observed in the discount parameters: 
individuals with lower institutional trust, lower education, or irregular preventive-health 
routines consistently demanded higher compensation for delay.

\paragraph{Economic Interpretation of MWTA and Heterogeneous Delay Sensitivity.}
MWTA provides a monetary expression of the behavioural cost of waiting and translates the
estimated delay coefficients into an economically interpretable metric. Subgroup comparisons
reveal pronounced heterogeneity. Respondents with low trust, lower education, chronic
conditions, or irregular preventive routines exhibit substantially higher MWTA values,
indicating that these groups perceive delay as more costly and place greater value on
immediacy. These patterns align with their steeper short-run discounting and elevated
$\kappa$ estimates.

This heterogeneity has direct implications for the cost--benefit evaluation of vaccination
strategies. Groups with high MWTA values stand to gain disproportionately from even modest
reductions in waiting time, suggesting that interventions such as queue-shortening, priority
scheduling, or rapid-access appointments generate especially large behavioural returns in
these populations. Conversely, broad, non-targeted convenience improvements may deliver
smaller marginal benefits among low-MWTA groups who are more tolerant of delay.

Subgroup-specific MWTA estimates also inform the efficient targeting of financial or
quasi-financial incentives. Because high-MWTA groups respond more strongly to reductions in
delay---or to modest, immediate rewards---allocating a greater share of incentives or
operational resources toward these populations may maximise uptake per yuan spent. Taken
together, this framework provides an economically grounded rationale for identifying where
behavioural interventions, convenience enhancements, or early-rollout policies can generate
the highest return on public-health investment.

\begin{table}[!t]
\centering
\caption{Marginal willingness-to-accept (MWTA) for one month of delay}
\label{tab:mwta_subgroups}
\begin{threeparttable}
\begin{tabular}{lccc}
\toprule
Group & MWTA & SE & Sig. \\
\midrule
Overall sample              & 52.1  &  7.4  & *** \\
Low institutional trust     & 81.3  & 10.9  & *** \\
High institutional trust    & 34.2  &  6.8  & *** \\
Irregular physical activity & 68.9  &  9.5  & *** \\
Regular physical activity   & 31.4  &  5.7  & *** \\
\bottomrule
\end{tabular}
\begin{tablenotes}[flushleft]
\footnotesize
\item Note: MWTA is expressed in RMB. SE denotes standard error.
\end{tablenotes}
\end{threeparttable}
\end{table}

MWTA offers a practical lens for policy design. Because reducing waiting time is often 
far cheaper and operationally easier than providing financial incentives, MWTA helps 
identify when logistical improvements may deliver greater uptake per yuan spent. For the 
average respondent, a one-month reduction in delay generates the same behavioural gain as 
roughly 50~RMB in incentives, with even larger returns among high-delay-aversion groups. 
This makes queue-shortening, extended clinic hours, or streamlined scheduling potentially 
more cost-effective tools than additional cash rewards.

To aid interpretation, Figure~\ref{fig:policy_tradeoff_menu} visualizes a ``Policy Trade-Off
Menu'' that translates these MWTA estimates into an intuitive comparison between time and
money. Panel~A standardizes the utility gain from reducing waiting time by one month, and
Panel~B shows the cash incentive required to offset that same delay for each subgroup.

\begin{figure}[!t]
    \centering
    \includegraphics[width=\linewidth]{The Policy Trade-Off Menu_ Monetary Equivalents of Reducing Vaccination Delay by One Month.png}
    \caption{Policy trade-off menu: monetary equivalents of reducing vaccination delay by one month. Panel~A standardizes the utility gain from reducing waiting time by one month at 1.0 (dashed reference line). Panel~B shows the cash incentive required to offset the same delay across subgroups (e.g., about 81~RMB for low-trust respondents versus about 31~RMB for respondents with regular preventive routines).}
    \label{fig:policy_tradeoff_menu}
\end{figure}

\clearpage

\section{Epidemiological Impact: Integrating Behavioral Parameters into SEIR Dynamics}

To translate individual delay aversion into population-level health outcomes, we integrate 
the recovered structural parameters into a behaviorally-augmented Susceptible-Exposed-Infectious-Removed 
(SEIR) compartmental model. Traditional models treat vaccination as an exogenous 
policy lever; we endogenize this rate by defining an \textit{effective vaccination rate} 
$v_{eff}(t)$ that reflects the interaction between logistical capacity and behavioral 
acceptance:
\[
v_{\text{eff}}(t) = v_{\text{cap}}(t) \cdot P(\text{delay}(t), C(t)),
\]
where $v_{\text{cap}}(t)$ captures delivery capacity (e.g., staffing, doses) and $P(\cdot)$ 
is the uptake probability derived from our Random Utility framework \citep{mcfadden2001}. 

Our simulation demonstrates that even modest delay aversion ($\kappa = 0.225$) can 
fundamentally alter outbreak trajectories. Because vaccination is most epidemiologically 
valuable during the early growth phase, behavioral delay shifts effective immunization 
later into the epidemic curve, allowing more infections to occur while $R_t$ is high. 
This behavioral friction can lead to a \textbf{5--10\% increase in the cumulative attack 
rate} compared to a counterfactual scenario where uptake is driven purely by logistical 
capacity. We acknowledge that our simulation assumes a constant uptake probability; in reality, delay aversion may be prevalence-elastic, with rising infection rates potentially collapsing present-oriented biases as the cost of waiting becomes more salient. By bridging micro-level behavioral parameters with macro-level transmission 
dynamics, this framework identifies ``behavioral delay'' as a biologically consequential 
bottleneck for public health preparedness and a key lever for health policy intervention.

\clearpage

\section{Discussion}\label{sec:discussion}
\doublespacing

This study examined how individuals value timely vaccination for COVID-like diseases (CLD) 
using a discrete choice experiment embedded within a behavioural-economic framework. 
Three central insights emerge. First, respondents exhibited short-run impatience, discounting delayed protection and requiring compensation for even modest waiting times. Second, these patterns align with established theories of present bias and 
quasi-hyperbolic discounting, indicating that vaccine delay reflects systematic intertemporal 
preferences rather than solely access constraints. Third, pronounced demographic, behavioural, 
and institutional gradients show that delay aversion is socially patterned and linked to broader 
determinants of preventive behaviour. While these absolute valuations may be upward-biased 
due to hypothetical decision-making, the underlying behavioural patterns remain stable across 
subgroups and align with international evidence. \textit{Although stated preferences may not perfectly match real incentive-consistent 
behaviour, extensive evidence shows that DCE-derived valuations of delay and convenience 
reliably predict uptake patterns in preventive-health contexts} \cite{Johnson2013,Hensher2015,Kong2025}.


\subsection{Behavioural Mechanisms Underlying Vaccine Delay}

The pattern of discounting observed here is consistent with evidence that vaccination decisions place substantial weight on short-term frictions relative to delayed health benefits \cite{vanderPol2001,Attema2013,Attema2018}, indicating particularly strong preferences for immediacy in vaccination decisions. Respondents placed disproportionate weight on short-term inconveniences and side-effect concerns relative to delayed and uncertain health benefits. The steep early decline in the hyperbolic curve likely captures the cumulative psychological burden of logistical frictions, which multi-country evidence identifies as key bottlenecks in pandemic vaccination \cite{Kong2025}. These tendencies intensify under uncertainty and perceived risk ambiguity \cite{Brewer2017}. The pattern is consistent with experimental evidence on near-term impatience \cite{Laibson1997,ODonoghue1999} and helps explain the ``strategic delay'' behaviour documented in prior campaigns \cite{Lievre2024,Tran2025}. Recent research underscores that monetary incentives might not always resolve these behavioral bottlenecks and could even risk crowding out intrinsic motivation in certain contexts \cite{Yamamura2026}.

\paragraph{Subjective Waiting Time and the Behavioral Delay Multiplier (BDM).}
To connect discounting behaviour to experiential perceptions of delay, we compute the \textit{Behavioral Delay Multiplier} (BDM). As derived in Appendix C9, the BDM is defined as the ratio between the perceived and objective duration of waiting, providing a behavioural measure of \textit{subjective waiting time}.

While the average BDM is 1.29, this inflation is even more pronounced among vulnerable subgroups: low-trust respondents and individuals with chronic conditions exhibit BDM values approaching 1.5, meaning that a 30-day delay is experienced as the equivalent of 45 days. We interpret this divergence between objective and subjective time as a \textit{psychological time tax} disproportionately borne by socially and medically vulnerable populations. Incorporating subjective waiting time into policy design clarifies why some groups respond strongly to even modest changes in waiting time and supports prioritising delay-reducing interventions for these populations.



\subsection{Biological Vulnerability and Intertemporal Health Investment}

Intertemporal preferences for vaccination timing are best understood not as abstract measures of impatience, but as responses shaped by biological vulnerability and exposure risk. When individuals face elevated baseline health risks, each day of delayed protection carries a greater expected biological cost. Steeper short-run discounting, in this light, reflects an adaptive valuation of immediacy under uncertainty---a rational response to risk rather than a departure from optimal health investment \parencite{Grossman1972, Attema2018}.

This interpretation finds support in the elevated delay sensitivity observed among respondents with chronic medical conditions. Chronic illness---whether cardiovascular, metabolic, or respiratory---compromises immune function, diminishes physiological reserves, and heightens susceptibility to severe infection outcomes. For these individuals, remaining unprotected is not merely inconvenient; it represents a period of amplified biological exposure. From a health-capital perspective, poorer baseline health steepens the risk gradient over time, making delayed vaccination disproportionately costly. The stronger preference for immediacy in this group thus reflects risk-weighted decision-making rather than generalized behavioural bias \parencite{Bok2024, Brewer2017}.

Age-related patterns reinforce this logic. Both younger and older adults exhibit steeper discounting than middle-aged respondents, but for distinct biological and behavioural reasons. Among older adults, immunosenescence---the progressive decline in immune competence with age---elevates both infection risk and disease severity, providing a clear physiological basis for urgency. For younger adults, steeper discounting may instead reflect greater uncertainty about future behavioural follow-through or competing life demands. In both cases, time preferences track perceived vulnerability across the life course rather than reflecting uniform impatience \parencite{Lawless2013, Wang2018}.

Institutional trust introduces a further layer of complexity. Lower trust heightens uncertainty about whether delayed vaccination will ultimately deliver its promised benefits---whether due to concerns about vaccine availability, quality, or institutional follow-through. Under such uncertainty, discounting future protection more steeply becomes a coherent response: if the probability of receiving effective protection diminishes, the rational weight placed on that future benefit declines accordingly. In this context, \"delay aversion\" may not be a pure psychological trait but a rational response to perceived supply unreliability. Trust, in this framework, functions as intertemporal capital---stabilising expectations about future health returns and thereby attenuating present bias \parencite{Brewer2017, Betsch2018}.

Taken together, these findings suggest that heterogeneity in vaccination timing preferences reflects adaptive responses to biological and institutional risk rather than idiosyncratic variation in patience. Viewing delay aversion through this lens clarifies why identical waiting periods impose unequal burdens across populations: for the chronically ill, the elderly, and those lacking institutional trust, each day of delay carries greater weight. This perspective underscores the importance of integrating biological vulnerability into analyses of intertemporal health behaviour and challenges interpretations that treat time preferences as stable, context-independent traits.

\subsection{The Unequal Burden of Delay: Quantifying the Regressive ``Time Tax''}

Waiting for vaccination is not costless. Our MWTA estimates quantify this burden: respondents required approximately 52~RMB per month to tolerate delay, a valuation comparable to modest cash incentives offered during actual vaccine campaigns. This figure captures something beyond monetary trade-offs---it reflects how individuals experience waiting as a salient friction, even when vaccines are nominally free.

Real-world benchmarks contextualise this magnitude. During the COVID-19 rollout, local governments in China and the United States deployed cash transfers of 200--300~RMB (USD~30--45) to encourage uptake \cite{Meng2021,Thirumurthy2022}. That our MWTA estimates fall below these thresholds suggests stated preferences may conservatively estimate the true experiential cost of delay, particularly in high-uncertainty environments. Even so, the results demonstrate that modest reductions in waiting time can meaningfully shift the perceived attractiveness of vaccination.

The aggregate, however, conceals pronounced inequality. As Figure~\ref{fig:policy_tradeoff_menu} illustrates, the value placed on immediacy is socially patterned. High-trust individuals value a one-month reduction in delay at approximately 34.2~RMB; low-trust individuals value the same reduction at 81.3~RMB---a 2.4-fold difference. Waiting, in effect, functions as a regressive ``Time Tax,'' extracting greater psychological cost from those who already face heightened uncertainty about future protection \cite{Betsch2018}.

We use ``Time Tax'' metaphorically rather than literally: the burden is experiential, not financial. Yet the implication is concrete. Identical delays generate vastly different consequences across groups---what constitutes a minor inconvenience for the trusting becomes a substantial barrier for the sceptical.

These patterns carry a clear message: delays in vaccination access are not neutral administrative features. They are uneven temporal burdens, shaped by vulnerability and institutional trust. Uniform rollout schedules that ignore this heterogeneity risk concentrating the costs of waiting precisely where they are hardest to bear---among populations already navigating greater uncertainty and fewer resources for timely protection.
\paragraph{Policy Actions}

These findings indicate that uniform vaccination rollout strategies are experienced very
differently across populations when supply or delivery capacity is constrained. Groups
characterized by high delay aversion—such as individuals with lower institutional trust or
lower educational attainment—exhibit particularly strong behavioral responses to reductions
in waiting time. For these populations, prolonged delays translate into disproportionately
large experiential and behavioral barriers to uptake. Institutional features that shorten
waiting times, including fast-tracking mechanisms, priority appointment slots, or mobile
vaccination units, may therefore substantially reduce barriers to timely protection among
delay-sensitive groups.

By contrast, individuals with lower sensitivity to delay—such as those with higher
institutional trust or more regular preventive health routines—appear more tolerant of
moderate waiting periods. For these groups, maintaining transparent, predictable, and
reliable scheduling may be sufficient to sustain vaccination uptake, even in the presence
of limited capacity. Recognizing this heterogeneity highlights how uniform delivery
schedules can concentrate the burdens of delay among populations already facing greater
uncertainty or vulnerability. Accounting for group-specific sensitivity to waiting time
clarifies how vaccination delivery arrangements shape unequal experiences of access and
timeliness across the population.

\subsection{Cost Sensitivity and Budget Constraints}

Sensitivity to monetary incentives ($\beta_{cash}$) is the denominator of our MWTA 
calculation. While our main model estimates an average effect, individuals from lower-income 
households may exhibit higher marginal utility for cash, leading to a lower implied MWTA 
relative to wealthier respondents. This implies that identical monetary incentives have 
varying ``purchasing power'' for timely protection across socioeconomic strata. This 
heterogeneity underscores the need for ``equity-weighted'' incentive schemes that account for 
budget constraints and ensure that vulnerable populations are not priced out of timely 
protection by higher delay aversion.

\subsection{Limitations}

Several limitations should be noted. First, as with all stated-preference studies, hypothetical bias may inflate absolute MWTA and discounting magnitudes. However, the survey design incorporated multiple features intended to approximate real-world vaccination decisions, including realistic attribute ranges informed by pilot testing, repeated choice scenarios, dominance checks, and an opt-out option. The very low rate of dominance violations and the consistency of subgroup gradients suggest that the behavioural patterns identified are robust, even if absolute valuations are upward-biased \cite{Louviere2000,Johnson2013}.

Second, the sample comprises adults in Wuhan, a high-salience pandemic context, and the DCE captures static, one-shot trade-offs rather than dynamic influences such as evolving risk perceptions or social learning \cite{Lievre2024,Tran2025}. While absolute effect sizes may vary in lower-salience settings, the structural relationships between delay aversion, institutional trust, and preventive behaviour are theoretically portable and empirically observable across populations. Future work should examine how intertemporal preferences evolve across outbreak phases.

Third, the study relies on self-reported chronic disease status rather than verified medical records or biomarkers. Future research could strengthen the biological interpretation by incorporating objective health measures---such as HbA1c for diabetes management or inflammatory markers---to validate the link between physiological vulnerability and delay aversion.


\section{Conclusion}

This study quantifies the intertemporal valuation of vaccination timing, treating delay not 
merely as a logistical friction but as a period of biological exposure. By recovering 
structural parameters ($\kappa = 0.225, \delta = 0.160$) and converting them into 
policy-relevant MWTA values, we demonstrate that individuals place a high ``shadow price'' 
on timely protection. These preferences are heterogeneous, regressive, and 
systematically patterned by biological vulnerability and institutional trust.

Our integrated behavioral-epidemiological framework highlights that neglecting delay 
aversion can lead to underestimation of outbreak risks and overestimation of vaccine 
rollout effectiveness. Policymakers should design ``time-consistent'' interventions that 
minimize the behavioral burden of waiting through prioritized scheduling, dynamic 
incentives, and clear risk communication targeted at high-delay-aversion populations. 
As the global public health community prepares for future infectious threats, 
incorporating the shadow price of delay will be essential for creating more resilient 
and equitable healthcare systems.





\section*{Declaration of Generative AI Use}

Generative artificial intelligence tools (including ChatGPT) were used to assist with language
editing, formatting, and clarification of exposition during manuscript preparation. These tools
were not used to generate data, conduct analyses, interpret results, or draw scientific
conclusions. All analyses and interpretations were performed by the authors, who take full
responsibility for the content of the manuscript.

\section*{Ethics approval and informed consent}

The study involved anonymous, minimal-risk survey research conducted in China. All study
procedures complied with relevant national regulations and institutional guidelines. The study
protocol was reviewed by an appropriate institutional ethics committee and determined to
require no formal ethics approval. All participants provided informed consent prior to
participation.

\section*{Data Availability Statement}

The data that support the findings of this study are not publicly available due to ethical and privacy considerations related to human-subject survey data. De-identified data may be made available from the corresponding author upon reasonable request, subject to approval by the relevant ethics committee.



\clearpage






\bibliographystyle{elsarticle-num}
\bibliography{references_cleaned}

\clearpage

\appendix
\appendix
\section*{Appendix A. DCE Instrument, Choice Tasks, and Validity Diagnostics}
\addcontentsline{toc}{section}{Appendix A. Internal Validity Diagnostics and DCE Materials}

This appendix presents the diagnostic checks used to assess internal validity in the discrete 
choice experiment (DCE), including (i) dominance behaviour, (ii) attribute non-attendance 
(ANA), (iii) transitivity and consistency, (iv) intra-individual variable correlations, and 
(v) pilot-study procedures. Collectively, these checks support the internal reliability and 
behavioural coherence of the responses.

\subsection*{A.1 Dominance Behaviour}

Dominance tests evaluate whether respondents ever selected an alternative that was strictly 
dominated across all attributes. As shown in Table~A1, only 1.8\% of respondents violated 
dominance at least once—well within the $\leq 5\%$ threshold commonly considered acceptable 
in DCE research \cite{Louviere2000,Johnson2013}. This low rate suggests that respondents 
understood the choice tasks and engaged in meaningful trade-offs.

\subsection*{A.2 Attribute Non-Attendance (ANA)}

ANA analyses assess whether respondents systematically ignored particular attributes during 
choice tasks. Table~A1 indicates that ANA was 0\% for vaccine delay, efficacy, side-effects, 
and cash incentives, reflecting strong salience of the primary decision attributes. A small 
minority (5.4\%) showed insensitivity to vaccine origin, suggesting minor simplification but 
no evidence of widespread attribute non-attendance. Overall, ANA diagnostics demonstrate 
robust engagement with the full attribute structure.

\begin{table}[H]
    \centering
    \caption{Dominance behaviour and attribute non-attendance (ANA)}
    \label{tab:validity}
    \begin{tabular}{lcc}
        \toprule
        \textbf{Check} & \textbf{Definition} & \textbf{Result (\%)} \\
        \midrule
        Dominance rate &
        Share selecting a dominated alternative &
        1.8 \\
        ANA: Vaccine delays &
        Choices invariant to delay levels &
        0.0 \\
        ANA: Vaccine efficacy &
        Choices invariant to efficacy levels &
        0.0 \\
        ANA: Side effects &
        Choices invariant to side-effect levels &
        0.0 \\
        ANA: Cash incentives &
        Choices invariant to cash levels &
        0.0 \\
        ANA: Vaccine origin &
        Choices invariant to origin levels &
        5.4 \\
        \bottomrule
    \end{tabular}

    \vspace{4pt}
    \footnotesize
    \textit{Note.} Dominance reflects selecting an option strictly inferior across all attributes. ANA 
    indicates choices unaffected by variation in a given attribute across tasks.
\end{table}

\subsection*{A.3 Transitivity and Choice Consistency}

Respondents’ choices satisfied the transitivity assumptions of random utility theory. No 
systematic violations or preference reversals were detected, and repeated patterns across the 
six tasks reflected stable and ordered preferences. These findings support the internal 
consistency and construct validity of the DCE, reinforcing the reliability of the resulting 
preference and time-discount parameter estimates.

\subsection*{A.4 Correlation Matrix of Intra-Individual Variables}

Figure~A1 displays pairwise correlations across all individual-level covariates included in 
the mixed-logit and subgroup models. All coefficients were below 0.10, indicating minimal 
concerns regarding multicollinearity or redundancy.

Access to healthcare showed only weak associations with other variables. Government trust was 
slightly positively correlated with age and rural residence, while smoking exhibited a weak 
negative correlation with chronic medical conditions. Overall, these low correlations confirm 
sufficient distinctness across covariates for use in the multivariate analyses.

\begin{figure}[H]
    \centering
    \includegraphics[width=0.9\textwidth]{correlation among Individual variables4.png}
    \caption{Correlation matrix of individual-level covariates. Darker colours indicate stronger 
    correlations. All coefficients are below 0.10, confirming negligible multicollinearity.}
    \label{fig:correlation}
\end{figure}

\subsection*{A.5 Example Choice Task}

Figure~A2 provides an example of the choice task presented to respondents. Each task 
featured two hypothetical vaccine profiles varying across five attributes—delay to 
vaccination, origin, efficacy, side-effect severity (supported by illustrative icons), and 
cash incentives—and an opt-out alternative. Attribute descriptions appeared at the top of the 
screen. For example, efficacy was defined behaviourally as the percentage reduction in 
infection risk (``95\% effectiveness means a 95\% reduction in infection risk''), and side 
effects were described using both text and images to reduce cognitive burden.

The example shown contrasts two vaccines identical in origin and efficacy but differing in 
delay (0 vs.\ 1 month), side-effect severity (fatigue for one week vs.\ fever lasting two 
days), and incentives (800 CNY vs.\ 200 CNY).

\begin{figure}[H]
\centering
\includegraphics[width=0.95\textwidth]{sample card.png}
\caption{Example DCE choice task shown to respondents. Participants selected between two hypothetical 
vaccine profiles or an opt-out option.}
\label{fig:sampletask}
\end{figure}

\subsection*{A.6 Pilot Study Procedures}

A pilot study was conducted with 60 adults from Wuhan to assess the clarity, feasibility, and 
cognitive burden of the DCE instrument. The pilot followed the same screening criteria as the 
main study and included the full sequence of choice tasks, sociodemographic questions, and a 
debriefing module.

The pilot served three purposes. First, it evaluated comprehension of attribute descriptions 
and level wording. Respondents demonstrated clear understanding of delay, efficacy, and 
side-effect attributes, while minor revisions were made to simplify phrasing for vaccine 
origin and severity icons.

Second, the pilot assessed task difficulty and respondent burden. Completion times ranged 
from $4.7$ to $10.2$ minutes (mean: $6.1$)
, indicating that the number of tasks and cognitive 
demands were appropriate. Dominance and ANA rates were consistent with those in the main 
sample, suggesting no evidence of satisficing or disengagement.

Third, the pilot refined the statistical efficiency of the experimental design. Preliminary 
parameter estimates informed updated priors in the efficient design algorithm, improving 
attribute balance and orthogonality in the final DCE.

Overall, pilot results confirmed that the survey instrument was clear, feasible, and capable 
of generating high-quality preference data.

\clearpage
\section*{Appendix B. Sample Representativeness and Census Benchmarking}
\label{app:B}
\doublespacing

This appendix evaluates the representativeness of the analytic sample ($N = 1{,}027$) by 
comparing key demographic characteristics with benchmarks from the 2020 Wuhan Census \cite{WuhanStats2021}.
Sampling of the quotes ensured the alignment in gender, age, and district of residence. As shown in 
Table~\ref{tab:census_compare_B1}, the sample closely approximates the demographic profile 
of Wuhan’s adult population. Minor differences—such as a slightly higher proportion of 
college-educated respondents—are typical of online survey panels and were further examined 
through robustness and weighting analyses.

\begin{table}[H]
\centering
\caption{Comparison of Analytic Sample with 2020 Wuhan Census Benchmarks}
\label{tab:census_compare_B1}
\begin{tabular}{p{5cm}cc}
\toprule
\textbf{Characteristic} & \textbf{Sample (\%)} & \textbf{Wuhan Census (\%)} \\
\midrule

\textit{Gender} & & \\
\quad Female & 53.0 & 51.1 \\
\quad Male   & 47.0 & 48.9 \\

\addlinespace
\textit{Age Distribution} & & \\
\quad 18--34 & 21.0 & 23.9 \\
\quad 35--44 & 31.2 & 28.6 \\
\quad 45--54 & 22.5 & 21.7 \\
\quad 55--65 & 13.2 & 14.3 \\
\quad 65+    & 12.0 & 11.5 \\

\addlinespace
\textit{Residence} & & \\
\quad Urban & 47.0 & 48.3 \\
\quad Rural & 53.0 & 51.7 \\

\addlinespace
\textit{Education Level} & & \\
\quad No formal schooling & 4.0 & 5.6 \\
\quad Primary school & 22.0 & 24.3 \\
\quad Middle school & 25.7 & 27.1 \\
\quad High school & 20.4 & 21.7 \\
\quad College/University & 15.6 & 13.4 \\
\quad Graduate degree & 11.0 & 7.9 \\

\bottomrule
\end{tabular}

\vspace{6pt}
\footnotesize
\textit{Note.} Census benchmarks are drawn from 
the Wuhan Bureau of Statistics\cite{WuhanStats2021}. Deviations follow typical patterns in online survey panels. 
Post-stratification weighting yielded mixed-logit and discount-rate estimates substantively 
identical to unweighted results (Appendix~C).
\end{table}

\begin{tcolorbox}[
    colback=gray!5,
    colframe=gray!50,
    boxrule=0.6pt,
    arc=3pt,
    left=6pt,
    right=6pt,
    top=6pt,
    bottom=6pt
]
\textbf{Summary of Sample Representativeness}

\begin{itemize}[leftmargin=1.2em]
    \item Survey sampling produced close alignment with Wuhan’s 2020 Census for gender, 
          age, and urban--rural residence.
    \item Slight upward shift in college-educated respondents is consistent with online 
          panel sampling.
    \item Pearson $\chi^2$ tests indicate no statistically significant difference between 
          sample and population distributions at the 5\% level.
    \item Weighted and unweighted mixed-logit models produced nearly identical results, 
          confirming robustness to demographic composition (see Appendix~C).
\end{itemize}
\end{tcolorbox}

To formally evaluate representativeness, Pearson $\chi^2$ tests were conducted for gender, 
age group, and residence. All tests were statistically non-significant at the 5\% level, 
indicating no systematic deviation from Wuhan’s adult population. Weighted mixed-logit 
models yielded parameters nearly identical to those in the unweighted sample, reinforcing 
the robustness of the main findings.


\clearpage
% ============================
% Appendix C. Model Estimation Details
% ============================

\section*{Appendix C Model Estimation Details}
\label{app:C}
\doublespacing

This appendix provides additional detail on the estimation of the mixed logit models and the
derivation of exponential and hyperbolic discount parameters used in the main analysis.

\subsection*{C1. Utility Specification}

We adopt a random–utility framework in which individual $i$ evaluates vaccine alternative $j$
in choice task $t$ according to:

\begin{align}
U_{ijt} &= \beta_0
+ \beta_1\,\text{VaccineDelay}_{\text{std},ijt}
+ \beta_2\,\text{VaccineOrigin}_{\text{std},ijt}
+ \beta_3\,\text{VaccineEfficacy}_{\text{std},ijt} \notag\\
&\quad + \beta_4\,\text{SideEffects}_{\text{std},ijt}
+ \beta_5\,\text{CashIncentives}_{\text{std},ijt}
+ \varepsilon_{ijt}, \tag{C1}
\end{align}

where $\varepsilon_{ijt}$ follows an i.i.d.\ Type I extreme-value distribution.  
For comparison, both conditional logit (MNL) and mixed logit (MXL) models were estimated,
but the MXL is the preferred specification because it allows random coefficients to capture
unobserved preference heterogeneity.

In vector notation:

\begin{equation}
U_{ijt}
= \beta_0
+ \sum_{k=1}^{K} \beta_k X_{k,ijt}
+ \varepsilon_{ijt},
\tag{C2}
\end{equation}

Standardization of the attribute coefficients facilitates comparison across attributes and inputs
directly into the discount-rate estimation procedure described below.

\subsection*{C2. Discount Function Specification}

Observed delay levels $t \in \{0,1,3,6\}$ months are mapped to standard exponential and
hyperbolic discounting frameworks:

\begin{align}
D_{\text{Exp}}(t;\hat{\delta}) &= e^{-\hat{\delta} t}, \qquad \hat{\delta} \ge 0, \tag{C3}\\[4pt]
D_{\text{Hyper}}(t;\hat{\kappa}) &= \frac{1}{1+\hat{\kappa} t}, \qquad \hat{\kappa} \ge 0. \tag{C4}
\end{align}


These formulations correspond to time-consistent (exponential) and present-biased
(hyperbolic) discounting, respectively.

\subsection*{C3. Constructing Observed Discount Points}

Let $\beta_{\text{delay (per month)}}$ be the unstandardized coefficient from the secondary model in
which delay is entered linearly. This coefficient represents the marginal utility loss associated
with a one-month increase in delay. It is recovered from the standardized coefficient via:

\[
\beta_{\text{delay (per month)}} 
= \frac{\beta_{\text{delay (std)}}}{\mathrm{SD}(\text{MultipleDelays})}. \tag{C5}
\]

Following standard practice in preventive-health time-preference studies, we define the
utility-consistent “observed” discount factor for each delay $t$ as:

\[
D_{\text{obs}}(t) 
= \exp\!\left(-\beta_{\text{delay (per month)}}\, t\right). \tag{C6}
\]

These points serve as empirical inputs for fitting exponential and hyperbolic discount functions.

\subsection*{C4. Estimation of Discount Parameters}

The exponential and hyperbolic discount parameters are obtained by minimizing the sum of
squared errors (SSE) between the observed and theoretical discount functions.

\paragraph{Hyperbolic fit.}

\[
\mathrm{SSE}_{\text{Hyper}}(\hat{\kappa})
= \sum_{i=1}^{n}
\left[
D_{\text{obs}}(t_i)
- D_{\text{Hyper}}(t_i;\hat{\kappa})
\right]^2 .
\tag{C7}
\]


The first-order condition is:

\[
\frac{\partial}{\partial \kappa}
\mathrm{SSE}_{\text{Hyper}}(\kappa)
= 2 \sum_{i=1}^{n}
\left(
D_{\text{obs}}(t_i)
- \frac{1}{1 + \kappa t_i}
\right)
\frac{t_i}{(1 + \kappa t_i)^2}
= 0,
\qquad \text{evaluated at } \kappa = \hat{\kappa}.
\tag{C8}
\]


\paragraph{Exponential fit}

\[
\mathrm{SSE}_{\text{Exp}}(\hat{\delta})
= \sum_{i=1}^{n}
\left[
D_{\text{obs}}(t_i)
- D_{\text{Exp}}(t_i;\hat{\delta})
\right]^2 .
\tag{C9}
\]

First-order condition:

\[
\frac{\partial}{\partial \delta}
\mathrm{SSE}_{\text{Exp}}(\delta)
= 2 \sum_{i=1}^{n}
\left(
D_{\text{obs}}(t_i)
- e^{-\delta t_i}
\right)
t_i e^{-\delta t_i}
= 0,
\qquad \text{evaluated at } \delta = \hat{\delta}.
\tag{C10}
\]

Closed-form solutions do not exist; parameters were estimated via nonlinear least squares.

\subsection*{C5. Subgroup Estimation and Multiplicity}

Subgroup-specific MXL models were estimated separately for demographic, behavioural,
health-status, and institutional-trust categories. Because this involves multiple hypothesis tests,
$p$-values for subgroup delay coefficients were adjusted using the Benjamini–Hochberg
false-discovery rate (FDR) at $q = 0.10$. Robustness of subgroup differences was confirmed
under Holm–Bonferroni corrections.


\subsection*{C.6 Recovering Discount Parameters from Waiting-Time Coefficients}

To recover exponential discount parameters from the DCE model, we map the estimated 
waiting-time coefficient $\beta_{\text{delay}}$ to the implied monthly discount factor $\delta$. 
Under exponential discounting, utility declines with delay according to $U(t) = U_0 \delta^{t}$. A first-order approximation yields a linear representation in which 
the marginal disutility of delay corresponds to the logarithm of the discount factor. 
Normalizing the baseline utility $U_0 = 1$ implies:

\[
\hat{\delta} = \exp\!\left(\hat{\beta}_{\text{delay}}\right).
\]


Thus, more negative values of $\beta_{\text{delay}}$ correspond to lower discount factors and 
stronger impatience toward delayed protection. This transformation is used to obtain the 
discount parameters reported in Table~\ref{tab:subgroups_long}.

\subsection*{C.7 MWTA Derivation}

The marginal willingness-to-accept (MWTA) for one month of delay is derived as the 
ratio of the delay coefficient to the cash-incentive coefficient:

\[
\text{MWTA} = -\frac{\beta_{\text{delay}}}{\beta_{\text{cash}}}.
\]

This represents the monetary compensation required to offset the disutility of an 
additional month of waiting.



\subsection*{C8. Additional Diagnostics}

We conducted additional diagnostic checks, including model comparison (AIC and
log-likelihood), scale-heterogeneity assessment using a GMNL specification, and
re-estimation excluding potential speeders. Across all alternative specifications, the sign,
magnitude, and statistical significance of the delay coefficient remained highly stable,
and the implied discount parameters were unchanged. These results confirm that the
main findings are robust to functional-form assumptions and response-quality concerns.
Full model outputs are available upon request.



\subsection*{C9. Behavioral Delay Multiplier (BDM)}

To provide an intuitive interpretation of the hyperbolic discount parameter, we compute a
Behavioral Delay Multiplier (BDM), defined as:

\[
\mathrm{BDM} = \frac{1}{1 - \hat{\kappa}}, \tag{C11}
\]

which rescales the estimated hyperbolic parameter into a measure of the perceived cost of
waiting. A value greater than 1 indicates that each month of objective delay is experienced as
a longer \emph{behavioural} delay. For example, the overall estimate $\hat{\kappa} = 0.225$
implies a BDM of 1.29, meaning that one month of waiting feels equivalent to approximately
1.29 months. Subgroups with larger $\hat{\kappa}$—such as older adults, respondents with
primary education, or those reporting neutral trust in government—exhibit BDM values
between 1.45 and 1.55, indicating substantially amplified perceived waiting costs. In contrast,
groups with smaller $\hat{\kappa}$, including rural residents and individuals reporting regular
physical activity, show BDM values near 1.15–1.20, reflecting weaker behavioural sensitivity
to delay. This multiplier offers a transparent summary measure that complements the formal
discount-rate estimates presented in Sections~C2–C4.

\begin{table}[H]
\centering
\caption{Behavioral Delay Multiplier (BDM) for Selected Subgroups}
\label{tab:BDM}
\begin{tabular}{lccc}
\toprule
\textbf{Subgroup} & $\hat{\kappa}$ & \textbf{BDM} & \textbf{Interpretation} \\
\midrule
Overall sample          & 0.225 & 1.29 & 1 month feels like 1.29 months \\
Urban residents         & 0.290 & 1.41 & Higher perceived waiting cost \\
Rural residents         & 0.174 & 1.21 & Lower perceived waiting cost \\
Primary education       & 0.348 & 1.53 & Amplified behavioural delay \\
Adults aged 65+         & 0.327 & 1.49 & Delay feels substantially longer \\
Neutral trust in gov.   & 0.356 & 1.55 & Strongest impatience \\
\bottomrule
\end{tabular}
\end{table}



\clearpage

% ===========================
% Appendix D: Subgroup-specific Discounting Plots
% ===========================
\section*{Appendix D. Subgroup-specific Discounting Plots}
\mbox{}
\newcommand{\panelpair}[4]{%
  \begin{minipage}{0.48\textwidth}
    \includegraphics[width=\linewidth]{#1}
    \caption*{#3}
  \end{minipage}\hfill
  \begin{minipage}{0.48\textwidth}
    \includegraphics[width=\linewidth]{#2}
    \caption*{#4}
  \end{minipage}%
}

% Hide subfigure labels (A, B, …) just for these figures
\captionsetup[subfigure]{labelformat=empty}

\begin{figure*}[t]
  \centering

  % Panels A–C
  \panelpair
    {Age Hyperbolic Discounting.png}
    {Age Exponential Discounting.png}
    {Panel A: Age-specific discounting. Five age groups (18–34, 35–44, 45–54, 55–65, 65+).}
    {Five age groups (18–34, 35–44, 45–54, 55–65, 65+).}

  \vspace{0.6em}
  \panelpair
    {Urban vs Rural-Hyperbolic discounting.png}
    {Urban vs Rural-Exponential discounting.png}
    {Panel B: Location-specific discounting}
    {Urban vs rural respondents.}

  \vspace{0.6em}

  \panelpair
    {Medical condition-Hyperbolic discounting.png} % Corrected filename
    {Medical condition-Exponential discounting.png}
    {Panel C: Medical condition–specific discounting}
    {Individuals with vs without medical conditions.}

  \caption{\textbf{Appendix Figure D1.} Subgroup-specific discounting (Panels A–C). Each panel compares $D(t)$ under hyperbolic vs exponential fits.}
  \label{fig:appendix_discounting_p1}
\end{figure*}
\FloatBarrier

\begin{figure*}[t!]
  \centering

  % Panels D–E
  \panelpair
    {Government trust-Hyperbolic discounting.png}
    {Government trust-Exponential discounting.png}
    {Panel D: Government trust–specific discounting}
    {Five levels of government trust.}

  \vspace{0.6em}

  \panelpair
    {Smoking status-Hyperbolic Discounting.png}
    {Smoking Status-Exponential Discounting.png}
    {Panel E1: Smoking status–specific discounting}
    {Never, sometimes, often.}

  \vspace{0.6em}

  \panelpair
    {Physical activity-Hyperbolic Discounting.png}
    {Physical activity-Exponential Discounting.png}
    {Panel E2: Physical activity–specific discounting}
    {Rare, sometimes, quite often.}

  \caption{\textbf{Appendix Figure D2.} Subgroup-specific discounting (Panels D–E). Each row shows hyperbolic vs exponential discount functions $D(t)$ for a category.}
  \label{fig:appendix_discounting_p2}
\end{figure*}

% Restore normal subfigure caption labels afterward (optional)
\captionsetup[subfigure]{labelformat=parens,labelsep=space}
\FloatBarrier

\clearpage
%
% ===========================
% Appendix E: Trade-off computation between waiting time and monetary incentive
% ===========================
\section*{Appendix E. Trade-off Computation between Waiting Time and Monetary Incentive}

The marginal rate of substitution (MRS) between waiting time (vaccine delay) and cash incentive
was computed as:
\begin{equation}
\text{MRS}_{\text{delay--cash}} = -\frac{\beta_{\text{cash}}}{\beta_{\text{delay}}}.
\end{equation}
This measures the change in waiting time equivalent to a 100~RMB change in cash incentive,
holding other vaccine attributes constant.

From the mixed logit estimates ($\beta_{\text{delay}}=-0.125$, $\beta_{\text{cash}}=0.102$),
we obtain:
\begin{equation}
\text{MRS}_{\text{delay--cash}} = -\frac{0.102}{-0.125} = 0.82.
\end{equation}

Hence, each 100~RMB incentive corresponds to a perceived reduction of 0.82 delay units.
Given that one delay unit in the DCE design is equivalent to approximately one month
($\approx115.6$~RMB), respondents value 115.6~RMB as equivalent to reducing vaccine delay by one month.
Confidence intervals were estimated using the delta method.


\end{document}

